\documentclass[11pt]{article}
\usepackage{amsmath,amssymb,amsthm,hyperref}
\newtheorem{theorem}{Theorem}
\newtheorem{lemma}[theorem]{Lemma}
\newtheorem{corollary}[theorem]{Corollary}
\begin{document}
\title{Priority Jets, Tail-Free Chain Extensions, and a Cover Obstruction}
\author{}
\date{November 24, 2026}
\maketitle

\paragraph{Status and scope.}
The identities and theorems below are complete written proofs, not
Lean-checked. They preserve the original admissible affine-ratio model.
They give a new three-head limiting compiler and an obstruction to a
specific proposed use of that compiler. They do not prove an improved
upper bound or a stronger lower bound for unrestricted head complexity.
The obstruction permits arbitrary adaptive priority orders and full
triangular supports, rather than merely rectangular supports.

\section{An exact two-head covariance jet}

Fix an order of the coordinates, initially the identity order, and write
$z_i=1-x_i$. Let $a_i,b_i$ be arbitrary real numbers, and let
$0<\epsilon<1$. Put
\[
 w_i=\epsilon^{i-1},\qquad
 D=\epsilon^n+\sum_{i=1}^n w_i z_i,\qquad
 A=\sum_{i=1}^n a_i w_i^2z_i,\qquad
 B=\sum_{i=1}^n b_i z_i,\qquad
 C=\sum_{i=1}^n a_i b_i w_i z_i.                 \tag{1}
\]
All four expressions are affine functions on the cube. Direct
subtraction gives the exact identity
\[
 J_{\epsilon,t}
 =\frac1t\left(\frac{A+tC}{D+tB}-\frac A D\right)
 =\frac{CD-AB}{D(D+tB)}.                         \tag{2}
\]
For fixed $\epsilon$, take $t>0$ sufficiently small that
$w_i+t b_i>0$ for every $i$. Both denominators then have strictly
negative slopes in the original variables $x_i$, and their minimum
on the entire cube is $\epsilon^n>0$. Thus (2) is the sum of two
actual admissible affine ratios. No coordinatewise-complement
symmetry of the model is assumed: the slopes have been checked
directly in the original variables.

As $t\downarrow0$, (2) converges uniformly on the finite cube to
$C/D-AB/D^2$. If the first missing coordinate is $i$, division by
$w_i$ and $w_i^2$, respectively, gives
\[
 \frac C D\longrightarrow a_i b_i,
 \qquad \frac A{D^2}\longrightarrow a_i
 \qquad(\epsilon\downarrow0).
\]
At the all-one input, $A=B=C=0$ for every parameter value.
Consequently the successive limits give the following exact limiting
function:
\[
 F(x)=
 \begin{cases}
 -a_i\displaystyle\sum_{j>i} b_j(1-x_j),
       &i\text{ is the first missing coordinate},\\
 0,    &x=(1,\ldots,1).
 \end{cases}                                      \tag{3}
\]
Given any desired positive error tolerance, first choose $\epsilon$
small enough for the second limit and then choose $t$ small enough
for the first limit and the strict slope inequalities. This realizes
that tolerance with two heads.

\paragraph{Two-missing-coordinate specialization.}
If the only missing coordinates are $i<j$, then $F=-a_i b_j$.
If one block precedes another, the mixed table therefore has rank one.
If the two blocks are interleaved, setting every $b_i=1$ and setting
$a_i=1$ on one block and $a_i=-1$ on the other gives a saturated
comparison of which missing coordinate appears first. This is a
forward comparison in deletion coordinates. It does not supply the
opposite orientation required by the existing reverse-grid gluing
theorem.

\section{Removing every unwanted coefficient}

Write $P_i=x_1\cdots x_i$, with $P_0=1$. Formula (3) is equivalently
\[
 F=-\sum_{1\le i<j\le n}a_i b_j
                      P_{i-1}(1-x_i)(1-x_j).       \tag{4}
\]
Set $a_0=0$. Expanding (4), the coefficient of the off-chain monomial
$P_i x_j$, where $0\le i\le n-2$ and $j\ge i+2$, is
\[
                         (a_{i+1}-a_i)b_j.          \tag{5}
\]
Indeed the positive contribution comes from the pair $(i+1,j)$,
and, when $i>0$, the negative contribution comes from $(i,j)$.
There are no other contributions. Every remaining monomial is one
of $P_0,\ldots,P_n$. For completeness, its coefficient is
\[
 c_0=-a_1\sum_{j>1}b_j,
 \qquad
 c_k=a_k\sum_{j>k}b_j-a_{k+1}\sum_{j>k+1}b_j
                  -\mathbf1_{k\ge2}a_{k-1}b_k
       \quad(1\le k\le n),                         \tag{6}
\]
where the expression involving $a_{n+1}$ is defined to be zero.

Any chain polynomial $\sum_{k=0}^n d_kP_k$ is a one-head limit.
Explicitly put $v_j=\sum_{k<j}d_k$ and use
\[
 \frac{\eta^n v_{n+1}+
             \sum_{j=1}^n\eta^{j-1}v_j(1-x_j)}
      {\eta^n+\sum_{j=1}^n\eta^{j-1}(1-x_j)}.        \tag{7}
\]
The denominator has strictly negative slopes and minimum $\eta^n$;
the ratio converges to the desired value at each first missing
coordinate and at the all-one input.

\begin{theorem}[Tail-free chain-extension compiler]\label{thm:compiler}
For arbitrary real arrays $u_0,\ldots,u_{n-2}$ and $b_1,\ldots,b_n$,
the polynomial
\[
 G(x)=\sum_{i=0}^{n-2}\sum_{j=i+2}^n
                                  u_i b_j P_i x_j                 \tag{8}
\]
is a uniform limit on the Boolean cube of sums of three admissible
heads. The same holds for every permutation of the coordinate order.
There are exactly the displayed coefficients and no other tails.
\end{theorem}

\begin{proof}
Choose $a_1=u_0$ and successively $a_{i+1}=a_i+u_i$.
By (5), the off-chain part of $F$ is exactly $G$. Subtract the
chain polynomial (6) using (7). Approximate $F$ by (2), and take
$\eta$ sufficiently small for the chain correction. The cube is
finite, so all errors are uniform. Relabeling the coordinates in
the displayed formulas preserves the checked common slope sign.
\end{proof}

As a check on indexing, direct exact Boolean inversion for arbitrary
integer arrays in dimensions $2$ through $7$ agrees with (5) and
(6). This finite calculation is not needed for the proof.

\section{A valid coefficient-to-value bridge}

There is a useful distinction between (8) and an earlier coefficient
kernel proposal. Here all unlisted coefficients have been removed
exactly in the limiting polynomial. Thus degree rescaling can be
implemented by resynthesizing the coefficient parameters, without
substituting non-Boolean inputs into a rational function.

\begin{lemma}[Compatible covers give strict signs]\label{lem:bridge}
Suppose a sum of $H$ polynomials of form (8) and $K$ chain
polynomials has a nonzero coefficient $c_S$ for every $S\subseteq[n]$.
Let $\sigma(S)=\operatorname{sgn}(c_S)$. Then a sum of at most
$3H+K$ actual admissible heads strictly has sign $\sigma(T)$ at
every Boolean input $\mathbf1_T$.
\end{lemma}

\begin{proof}
For $M>1$ multiply every coefficient of degree $k$ by $M^k$.
In (8), this is accomplished by replacing $u_i$ by $M^{i+1}u_i$;
chain coefficients can be scaled individually in (7). Thus the
rescaled polynomial is still in the same limiting family.
Let $a=\min_S|c_S|>0$ and $b=\max_S|c_S|$. At $\mathbf1_T$ its
value is $c_TM^{|T|}+\sum_{S\subsetneq T}c_SM^{|S|}$. For nonempty
$T$, the absolute value of the latter sum is at most
$2^n b M^{|T|-1}$. Choose $M>2^n b/a$. The leading term then
strictly determines the sign; the empty input has value $c_\varnothing$.
There is a positive minimum sign margin on the finite cube.
Theorem~\ref{thm:compiler} and (7), with total approximation error
smaller than that margin, give the claimed actual heads.
\end{proof}

In particular a cover of the monomials by modules whose nonzero
coefficients all agree with a prescribed sign table would give
such a representation. The next section shows that this
noncancelling cover strategy cannot supply an extra logarithm,
even using the entire triangular support of (8).

\section{An adversarial sign table for all adaptive orders}

Call a subset of entries of an $n$ by $n$ matrix an ordered-prefix
support if a subset $I$ of the rows is chosen, kept in its original
order, and each column contains a prefix of those selected rows.
Different columns may have different prefix lengths, including zero.
Call its signs compatible if, on all selected entries, they have
the form $\alpha_i\beta_j$ for row and column signs.

\begin{lemma}\label{lem:random}
For every positive integer $n$, there exists an $n$ by $n$ sign
matrix $h$ for which every compatible ordered-prefix support
contains fewer than $8n$ entries.
\end{lemma}

\begin{proof}
Choose all entries of $h$ independently and uniformly from
$\{-1,1\}$. Set $m=8n$. There are at most $2^n$ choices of the
selected row set and at most $2^{2n}$ choices of the row and column
signs, using arbitrary extensions of the row signs outside $I$.
A support with $m$ entries is specified by nonnegative prefix
lengths $r_1,\ldots,r_n$ with sum $m$, of which there are
$\binom{m+n-1}{n-1}$. This overcounts invalid lengths exceeding
$|I|$, which is harmless. For each valid fixed support and signs,
the probability of agreement is $2^{-m}$, since the constrained
matrix entries are distinct independent signs. The union bound
is therefore at most
\[
 2^{3n}\binom{9n-1}{n-1}2^{-8n}
 \le \left(\frac{8\cdot9e}{256}\right)^n
 < \left(\frac{216}{256}\right)^n <1,             \tag{9}
\]
using $\binom{9n}{n}\le(9e)^n$ and $e<3$.
There is consequently an $h$ with no such support of exactly $8n$
entries. A compatible larger support can have its column prefix
lengths reduced until their sum is $8n$, while retaining
compatibility. Thus it also cannot exist.
\end{proof}

The prefix encoding in this proof is essential. A union bound
over the $n!$ priority permutations loses a logarithm unnecessarily.
When there are only $m$ constrained entries, the prefix lengths
encode the relevant support in $\binom{m+n-1}{n-1}$ possibilities.

Fix a matrix $h$ from Lemma~\ref{lem:random} and define a sign
table on all subsets by
\[
 \sigma(\varnothing)=1,\qquad
 \sigma(S)=\prod_{j\in S}h_{|S|,j}
                   \quad(S\ne\varnothing).        \tag{10}
\]

\begin{theorem}[Noncancelling chain-extension cover obstruction]
Every polynomial of form (8), in any coordinate order, whose
nonzero coefficient signs agree with (10), has fewer than $8n$
nonzero coefficients. Hence a coefficientwise sign-compatible
cover of (10) using $H$ such modules and $K$ arbitrary chain
polynomials must satisfy
\[
                         8nH+(n+1)K\ge2^n.          \tag{11}
\]
\end{theorem}

\begin{proof}
Let the order be $\pi$, and write
$P_i=\{\pi(1),\ldots,\pi(i)\}$ in set notation. The off-chain
support consists of sets $P_i\cup\{j\}$ with
$\operatorname{pos}_\pi(j)>i+1$. By (10),
\[
 \sigma(P_i\cup\{j\})
 =\left(\prod_{\ell\in P_i}h_{i+1,\ell}\right)h_{i+1,j}.
                                                        \tag{12}
\]
Select row $i+1$ precisely when $u_i\ne0$, and select a column
precisely when $b_j\ne0$. For any fixed column $j$, the selected
rows appearing in its support are exactly those with
$i+1<\operatorname{pos}_\pi(j)$, a prefix of the selected rows.
Agreement of the coefficient signs $\operatorname{sgn}(u_i b_j)$
with (12) says that the entries of $h$ on this support equal a
product of row and column signs. Lemma~\ref{lem:random} bounds its
size by $8n$. The mapping from pairs $(i,j)$ to monomials is
injective: the degree fixes $i$, and then $j$ is the unique
coordinate outside $P_i$. Finally each chain polynomial supports
at most $n+1$ coefficients. Covering all $2^n$ monomials yields
(11), even when compatible supports overlap.
\end{proof}

\paragraph{Exact remaining scope.}
The compiler itself is valid and has only three heads per module.
The obstruction rules out the proposed strategy that directly covers
the desired Boolean sign table by coefficientwise compatible modules
and then applies Lemma~\ref{lem:bridge}. It does not rule out signed
cancellations between modules, choosing a different coefficient sign
table that nevertheless represents the desired output signs, higher
priority jets, or general finite-scale affine ratios. Any of those
would be a materially different global construction. In particular,
(11) is not a lower bound on unrestricted $H^*$ or on the minimum
number of arbitrary chain-extension modules in a sign representation.

\end{document}
